\section{Vista de Atributos de Calidad}\label{sec:vista-calidad-de-servicio}

La Vista de Atributos de Calidad se centra en los requerimientos de atributos de calidad de la plataforma. Los atributos de calidad son propiedades medibles o comprobables de un sistema, que indican qué tan bien el sistema satisface las necesidades de las partes interesadas \cite{Bass2012}. Los requerimientos de esta sección surgen de la letra del laboratorio \cite{LINS2026} y se verifican mediante pruebas automatizadas, la prueba de performance y los mecanismos de observabilidad.

\begin{longtable}{|p{1.3cm}|p{2.7cm}|p{11.5cm}|}
    \caption{Atributos de Calidad}\label{tab:calidad}\\
    \hline
    \textbf{Id} & \textbf{Tipo} & \textbf{Descripción} \\
    \hline
    \endfirsthead
    \hline
    \textbf{Id} & \textbf{Tipo} & \textbf{Descripción} \\
    \hline
    \endhead
    \multicolumn{3}{|l|}{\textit{Seguridad}} \\
    \hline
    AC01 & Seguridad & Las contraseñas de los usuarios del backoffice las guarda el directorio LDAP con hash y \textit{salt} (R7); el componente central no las almacena. El directorio solo se accede por LDAPS. \\
    \hline
    AC02 & Seguridad & Las comunicaciones con los componentes periféricos, las plataformas externas y el componente móvil utilizan HTTPS (R8). \\
    \hline
    AC03 & Seguridad & La autorización se controla en la capa de servicios, con permisos asociados a roles (R9). \\
    \hline
    AC04 & Seguridad & La pista de auditoría permite reconstruir el usuario, la acción, el fundamento y la marca de tiempo de cada acción de fiscalización, y no puede modificarse ni eliminarse desde la aplicación (R10). Cada registro guarda el caso, la acción (crear, notificar, presentar descargo, confirmar o desestimar), el usuario y su rol, los estados anterior y nuevo, el fundamento y la fecha. Un disparador de PostgreSQL rechaza cualquier modificación o borrado de esos registros, aunque se intente desde fuera de la aplicación, y la pista referencia al caso solo por su identificador, para que borrar un caso no la borre en cascada. \\
    \hline
    AC05 & Seguridad & Las credenciales, claves y certificados se administran fuera del repositorio de código y de las imágenes de contenedor (R11). \\
    \hline
    AC06 & Privacidad & La información pública se publica agregada, respetando la normativa de protección de datos personales, con un umbral mínimo de agregación de \pendiente{definir y justificar el umbral (R12)}. \\
    \hline
    \multicolumn{3}{|l|}{\textit{Escalabilidad y performance}} \\
    \hline
    AC07 & Escalabilidad & El componente central escala horizontalmente en la capa web o de aplicaciones, sin estado de sesión almacenado localmente en las instancias. \\
    \hline
    AC08 & Performance & La plataforma sostiene, sin pérdida de eventos, un volumen de ingesta de tracking de \pendiente{estimar y justificar: cantidad de vehículos reportando en forma concurrente y frecuencia de reporte}, y picos de al menos tres veces ese volumen. \\
    \hline
    AC09 & Performance & El percentil 95 del tiempo de respuesta de las consultas del frontoffice es menor a 500 ms. \\
    \hline
    AC10 & Performance & Los tiempos de respuesta no se degradan a lo largo de la prueba de performance. \\
    \hline
    AC11 & Performance & La prueba de performance determina el punto de quiebre e identifica los cuellos de botella con evidencia de los mecanismos de observabilidad. \\
    \hline
    \multicolumn{3}{|l|}{\textit{Resiliencia e idempotencia}} \\
    \hline
    AC12 & Resiliencia & La indisponibilidad de un sistema externo o de un nodo periférico no provoca la indisponibilidad del componente central. \\
    \hline
    AC13 & Resiliencia & Las integraciones con sistemas externos y nodos periféricos aplican un tiempo de espera, reintentos y un corte ante fallas sostenidas, con un comportamiento degradado definido:
    \begin{itemize}[leftmargin=*,topsep=2pt,itemsep=1pt,parsep=0pt,partopsep=0pt]
        \item \textbf{Balanzas:} espera de 5 s, 2 reintentos a 1 s y 2 s, y corte por 60 s tras 5 fallas en 10 llamadas. La evaluación de exceso de carga queda pendiente y el resto de las reglas se evalúa igual.
        \item \textbf{PDI:} espera de 3 s, 1 reintento a 1 s, y corte por 30 s tras 5 fallas en 10 llamadas. El perfil del ciudadano se crea sin los datos de la PDI, que se completan en su próximo ingreso.
        \item \textbf{Usuario gub.uy:} espera de 5 s, sin reintentos. El ingreso falla con un mensaje y las sesiones ya iniciadas siguen funcionando.
        \item \textbf{Correo:} espera de 5 s y 3 reintentos a 1, 5 y 30 min. El caso avanza igual y la notificación se reintenta.
        \item \textbf{Servicio de rutas:} \pendiente{definir la espera, los reintentos y el corte}. La Guía se registra igual y el recorrido esperado queda pendiente hasta que se pueda obtener (CU-07).
        \item \textbf{Directorio LDAP:} \pendiente{definir la espera, los reintentos y el corte}. El ingreso al backoffice falla con un mensaje; las sesiones ya iniciadas siguen funcionando, porque el JWT no depende del directorio (DA03).
    \end{itemize}
    Los tiempos de espera, los reintentos y el corte de balanzas, PDI, Usuario gub.uy, el servicio de rutas y el directorio LDAP se implementan con la biblioteca de tolerancia a fallas que trae WildFly (MicroProfile Fault Tolerance), mediante anotaciones en los adaptadores (\texttt{@Timeout}, \texttt{@Retry}, \texttt{@CircuitBreaker}), y el comportamiento degradado con \texttt{@Fallback}. Los reintentos del correo, que se espacian en minutos, los realiza un temporizador que reprocesa los envíos pendientes. \\
    \hline
    AC14 & Resiliencia & Las pruebas incluyen un escenario con el Sistema de Balanzas indisponible y otro con la PDI sin responder. \\
    \hline
    AC15 & Idempotencia & Un evento recibido más de una vez no altera dos veces el estado del viaje ni genera casos de fiscalización adicionales. \\
    \hline
    AC16 & Consistencia & El orden de la ruta se reconstruye a partir del momento de generación de cada evento y no del orden de llegada. \\
    \hline
    AC17 & Resiliencia & Ante una saturación, las posiciones de tracking esperan en la cola y se procesan más tarde, sin perder ninguna. Para eso se usa el \emph{paging} del broker de mensajería (DA02, política \texttt{PAGE}): cuando la cola \texttt{queue\_tracking} llena su memoria, el broker guarda las posiciones nuevas en disco en lugar de descartarlas, y las procesa en orden a medida que se libera memoria. Se descartan tirar posiciones, porque la plataforma no debe perder eventos (AC08), y pedirle al nodo que reporte menos seguido, porque la frecuencia de reporte la fija el sistema de tracking, que es externo a la plataforma. \\
    \hline
    \multicolumn{3}{|l|}{\textit{Operación sin conectividad}} \\
    \hline
    AC18 & Disponibilidad & La aplicación móvil es utilizable sin conexión: conserva localmente la Guía descargada y los eventos pendientes. \\
    \hline
    AC19 & Idempotencia & Al recuperar la conectividad, la sincronización es idempotente: cada evento lleva un identificador generado en el dispositivo y su reenvío no produce duplicados. \\
    \hline
    AC20 & Consistencia & Ante un cambio hecho desde el frontoffice mientras el dispositivo estaba sin conexión, prevalece el componente central. El frontoffice solo modifica o elimina viajes que todavía no comenzaron, así que un evento sobre un viaje eliminado o reasignado queda en conflicto y no se aplica. Cada evento se procesa por separado y se registra con su resultado, sin pérdida de información, y la aplicación adopta el estado del viaje que devuelve el componente central (CU-09, Sección \ref{sec:cu-09}). Prevalece el componente central porque la asignación de los viajes la decide el responsable de la empresa. \\
    \hline
    \multicolumn{3}{|l|}{\textit{Observabilidad}} \\
    \hline
    AC21 & Observabilidad & Los componentes emiten logs estructurados que incluyen un identificador de correlación. El componente central y el nodo de tracking mandan sus logs con OpenTelemetry (OTLP) a la Plataforma de Observabilidad, como registros estructurados que llevan el identificador de la traza en curso: es el identificador de correlación, el mismo en el nodo y en el componente central para una misma posición. \\
    \hline
    AC22 & Observabilidad & El componente central expone métricas de eventos procesados, latencia de procesamiento, tamaño de las colas y errores por tipo, y endpoints de verificación de estado. Las métricas se exportan por OTLP con Micrometer: posiciones procesadas por resultado, errores por tipo, tiempo de procesamiento, espera en la cola y antigüedad de cada posición, mensajes y consumidores de la cola de tracking, y tiempo de respuesta de la API REST, con histogramas para calcular percentiles. Los endpoints de verificación de estado son \texttt{/health/live} y \texttt{/health/ready} del servidor de aplicaciones y \texttt{/api/salud}, que verifica PostgreSQL y el Almacén de Posiciones; la Plataforma de Observabilidad los consulta periódicamente. \\
    \hline
    AC23 & Observabilidad & Al menos un flujo que atraviesa el componente móvil, el componente central y un nodo periférico se traza de forma distribuida. Ese flujo es el fin de un viaje: el chofer lo registra en el móvil, el componente central lo recibe al sincronizar y el Motor de Fiscalización consulta los pesajes al nodo de Balanzas (CU-11, DA10). Además se traza la ingesta: el nodo de tracking abre una traza por cada lote que publica y su contexto viaja en cada mensaje, así que el procesamiento de cada posición en el componente central forma parte de esa traza. \\
    \hline
    AC24 & Observabilidad & El informe de la prueba de performance se apoya en la evidencia obtenida de logs, métricas y trazas. La plataforma incluye un tablero con esa evidencia para la ingesta, el nodo de tracking y la API REST. \\
    \hline
    \multicolumn{3}{|l|}{\textit{Pruebas y calidad del código}} \\
    \hline
    AC25 & Mantenibilidad & Las pruebas automatizadas cubren aproximadamente el 80~\% de la lógica, con al menos una prueba por servicio. La cobertura se mide con JaCoCo en el pipeline de integración continua. \\
    \hline
    AC26 & Corrección & Las pruebas cubren el ciclo de vida de los casos de fiscalización, verifican que no sean posibles transiciones de estado no previstas y comprueban la idempotencia del procesamiento de tracking y de la sincronización móvil. \\
    \hline
    AC27 & Mantenibilidad & El pipeline de integración continua ejecuta un análisis estático de calidad de código. \\
    \hline
\end{longtable}
